\PDEIndexMark{idx:riesz-representation-introduction}Hilbert 空间的对偶空间十分简单: 它与该空间本身等距. 这一结果称为 Riesz 表示定理, 是 Hilbert 空间理论的基础之一.

\begin{Theorem}{Riesz 表示}
\label{thm:hilbert-riesz-representation}
Hilbert 空间的对偶空间与该 Hilbert 空间本身等距. 特别地, 对每个 $x\in H$, 由下式定义的 $H$ 上的线性泛函
\begin{equation}\label{eq:hilbert-riesz-inner-product-functional}
l_x(y):=(x,y)
\end{equation}
有界, 且范数为 $\|l_x\|_{H^*}=\|x\|_H$. 此外, 对每个 $l\in H^*$, 存在唯一的 $x\in H$, 使
\begin{equation}\label{eq:hilbert-riesz-functional-representation}
l(y)=(x,y)\quad\text{for every }y\in H,
\end{equation}
并且有 $\|x\|_H=\|l\|_{H^*}$.
\end{Theorem}

\begin{Proof}
如定理的陈述所示, 对每个 $x\in H$, 可以利用 \eqref{eq:hilbert-riesz-inner-product-functional} 定义线性泛函 $l_x$. 由 Cauchy-Schwarz 不等式, 有 $l_x\in H^*$ 且 $\|l_x\|\leq\|x\|$. 此外, 令 $y=x$, 得到 $\|l_x\|=\|x\|$.

另一方面, 设 $l$ 为 $H$ 上的非零线性泛函. 定义 $M=\{y\in H\mid l(y)=0\}$. 那么 $M$ 为 $H$ 的闭子空间. 此外, $M^\perp$ 为一维, 因为对 $M^\perp$ 中任意 $u$ 和 $v$, 有 $l(v)u-l(u)v\in M$. (这意味着 $l(v)u-l(u)v=0$.)

现在设 $\widetilde x$ 为 $M^\perp$ 中的单位向量. 对每个 $y\in H$, 有分解 $y=(\widetilde x,y)\widetilde x+z$, 其中 $z\in M$, 因而 $l(y)=(\widetilde x,y)l(\widetilde x)$. 于是, 若令
\begin{equation}\label{eq:riesz-representing-vector-from-kernel-complement}
x=\overline{l(\widetilde x)}\widetilde x,
\end{equation}
就可看到 $l=l_x:=(x,\cdot)$. 范数的相等性已经得到证明.
\end{Proof}

\begin{Remark}
\label{rem:riesz-mappings}
在前面的证明中, 构造了 Hilbert 空间及其对偶的相应元素之间的同构映射 $A_H:H^*\to H$ 和 $A_H^{-1}:H\to H^*$; 它们称为 \emph{Riesz 映射}.
\end{Remark}

\Needspace{4\baselineskip}
\begin{Example}
\label{ex:canonical-bidual-isometry}
设 $X$ 为 Banach 空间. 那么每个 $x\in X$ 都通过对应关系 $l_xy=y(x)$ 定义了 $X^*$ 上的线性泛函. 显然 $\|l_x\|\leq\|x\|$, 并且由将在下一小节~\ref{subsec:hahn-banach-theorem} 陈述的 Hahn-Banach 定理 (Theorem~\ref{thm:hahn-banach-extension}), 存在 $y\in X^*$, 使 $\|y\|=1$ 且 $y(x)=\|x\|$. 这意味着 $\|l_x\|=\|x\|$, 因而得到从 $X$ 到 $X^{**}$ 的一个子空间的等距映射. 刻画这一等距映射为满射的那些 Banach 空间十分重要.
\end{Example}

\begin{Definition}
\label{def:reflexive-banach-space}
若由 $l_xy=y(x)$ 定义的从 $X\to X^{**}$ 的等距映射 $x\mapsto l_x$ 为满射, 则称 Banach 空间\textbf{自反}.
\end{Definition}

由 Theorem~\ref{thm:hilbert-riesz-representation} 显然可知, 所有 Hilbert 空间都自反. 此外, 空间 $L^p(\Omega)$, $1<p<\infty$, 也自反. 相比之下, 可以证明 $L^1(\Omega)$ 和 $L^\infty(\Omega)$ 不自反. 自反空间具有许多有用的性质, 因而便于分析. 在本书其余部分中, 将主要在 Hilbert 空间中工作, 不再进一步研究自反 Banach 空间. 请读者参阅有关文献, 例如 \cite{DS}.
